\documentclass[12pt]{article}
\usepackage[utf8]{inputenc}
\usepackage[T1]{fontenc}
\usepackage[francais]{babel}
\usepackage[margin=2cm]{geometry}
\usepackage{amsmath, amssymb}
\usepackage{tikz}
\usetikzlibrary{trees}

\title{Correction du TD 3 – Analyse Syntaxique}
\author{}
\date{}

\begin{document}

\maketitle

\section*{Exercice 1 : Comprendre les grammaires}

\subsection*{1. Exemples d'expressions valides}
\begin{itemize}
    \item \(a + b\)
    \item \(a * b + c\)
    \item \((a + b) * c\)
\end{itemize}

\subsection*{2. Terminaux et non-terminaux}
\begin{itemize}
    \item \textbf{Terminaux} : \(+, *, (, ), \text{id}\) (où \(\text{id}\) représente un identifiant comme \(a, b, c\)).
    \item \textbf{Non-terminaux} : \(\text{Expr}, \text{Term}, \text{Factor}\).
\end{itemize}

\subsection*{3. Symbole de départ}
Le symbole de départ est \(\text{Expr}\).

\section*{Exercice 2 : Arbre de dérivation}

\subsection*{1. Arbre syntaxique concret}
\begin{tikzpicture}[level distance=1.5cm, level 1/.style={sibling distance=3cm}, level 2/.style={sibling distance=1.5cm}]
    \node {Expr}
        child {node {Expr}
            child {node {Term}
                child {node {Factor}
                    child {node {\(a\)}}
                }
            }
        }
        child {node {\(+\)}}
        child {node {Term}
            child {node {Term}
                child {node {Factor}
                    child {node {\(b\)}}
                }
            }
            child {node {\(*\)}}
            child {node {Factor}
                child {node {\(c\)}}
            }
        };
\end{tikzpicture}

\subsection*{2. Arbre syntaxique abstrait (AST)}
\begin{tikzpicture}[level distance=1.5cm, level 1/.style={sibling distance=3cm}, level 2/.style={sibling distance=1.5cm}]
    \node {\(+\)}
        child {node {\(a\)}}
        child {node {\(*\)}
            child {node {\(b\)}}
            child {node {\(c\)}}
        };
\end{tikzpicture}

\subsection*{3. Priorité de la multiplication}
La multiplication est prioritaire car la grammaire impose que les termes (\(\text{Term}\)) soient évalués avant les expressions (\(\text{Expr}\)). Ainsi, \(b * c\) est regroupé avant l'addition \(a + \ldots\).

\section*{Exercice 3 : Ambiguïté de grammaire}

\subsection*{1. Exemple d'instruction ambiguë}
\[
\text{if } E \text{ then if } F \text{ then } A \text{ else } B
\]

\subsection*{2. Deux arbres de dérivation possibles}
\begin{itemize}
    \item \textbf{Interprétation 1} : Le \texttt{else} se rapporte au deuxième \texttt{if} :
    \begin{verbatim}
    if E then (if F then A else B)
    \end{verbatim}
    
    \item \textbf{Interprétation 2} : Le \texttt{else} se rapporte au premier \texttt{if} :
    \begin{verbatim}
    if E then (if F then A) else B
    \end{verbatim}
\end{itemize}

\subsection*{3. Règle pour lever l'ambiguïté}
La règle habituelle est d'associer le \texttt{else} au \texttt{if} le plus proche (\textit{closest nested rule}).

\section*{Exercice 4 : Grammaire non ambiguë et factorisée}

\subsection*{Grammaire modifiée}
\[
\begin{aligned}
\text{Expr} &\rightarrow \text{Term } \text{Expr}' \\
\text{Expr}' &\rightarrow + \text{Term } \text{Expr}' \mid \epsilon \\
\text{Term} &\rightarrow \text{Factor } \text{Term}' \\
\text{Term}' &\rightarrow * \text{Factor } \text{Term}' \mid \epsilon \\
\text{Factor} &\rightarrow (\text{Expr}) \mid \text{id}
\end{aligned}
\]

\section*{Exercice 5 : Analyse descendante vs ascendante}

\subsection*{1. Analyse descendante}
L'analyse descendante part du symbole de départ et applique les règles de production pour dériver la chaîne d'entrée. Elle est souvent utilisée pour les grammaires LL.

\subsection*{2. Analyse ascendante}
L'analyse ascendante part de la chaîne d'entrée et tente de la réduire au symbole de départ en utilisant les règles de production à l'envers. Elle est adaptée aux grammaires LR.

\subsection*{3. Avantages}
\begin{itemize}
    \item \textbf{Descendante} : Simple à implémenter pour les grammaires non ambiguës.
    \item \textbf{Ascendante} : Plus puissante, capable de gérer des grammaires complexes (ex. : LR(1)).
\end{itemize}

\section*{Exercice 6 (Bonus) : Analyse descendante}

\subsection*{1. Vérification de \(x y z\)}
La chaîne \(x y z\) est valide car elle peut être dérivée à partir de \(S\) :
\[
S \rightarrow A B \rightarrow x A B \rightarrow x y B \rightarrow x y z
\]

\subsection*{2. Tableau d'analyse}
\begin{tabular}{|c|c|c|c|}
\hline
Étape & Pile & Texte à analyser & Action \\
\hline
0 & \(S\) & \(x y z\) & \(S \rightarrow A B\) \\
1 & \(A B\) & \(x y z\) & \(A \rightarrow x A\) \\
2 & \(x A B\) & \(x y z\) & Consommer \(x\) \\
3 & \(A B\) & \(y z\) & \(A \rightarrow y\) \\
4 & \(y B\) & \(y z\) & Consommer \(y\) \\
5 & \(B\) & \(z\) & \(B \rightarrow z\) \\
6 & \(z\) & \(z\) & Consommer \(z\) \\
7 & (vide) & (vide) & Acceptation \\
\hline
\end{tabular}

\subsection*{3. Conclusion}
La chaîne \(x y z\) est valide car l'analyse descendante aboutit à une acceptation complète.

\end{document}
